\part{Exodus}

\chapter{No Sympathy}

Ein Junge im schwarzen T-Shirt, etwa drei Ellen groß, mit schwarzen Haaren und pechschwarzen Iriden unter einer Sonnenbrille rückte sich die schwarzen Socken unter der schwarzen Jeans zurecht, damit diese besser in die schwarzen Klettschuhe passten. Dann erhob er sich und blickte sich um. Seine Kleidung schluckte das einfallende Licht und wäre selbst im Dunkeln durch starken Kontrast aufgefallen. So anziehend wie das alles für das Licht wirkte, so abstoßend war es für die Mitmenschen im Restaurant. Fnbsrd hatte sich bisher bei jeder Gelegenheit äußerst unbeliebt gemacht; nur die Zivilisiertheit der fortschrittlichen Stadt hatte einen Lynchmord verhindert.

Der ältere Herr am Tisch zu seiner Rechten sah einigermaßen sympathisch aus.

»Verzeihen Sie. Stört es Sie, falls ich mich zu Ihnen setze?«

»Wenn jemand fragt, ob sein Gegenüber eine ehrliche oder eine freundliche Antwort erwartet, dann steckt im Wort ›oder‹ bereits eine unfreundliche Antwort.«

Das klang sehr allgemein gehalten. »M… hm?«

»Und deshalb kann ich das nicht fragen.«

Es handelte sich um die wohl freundlichste Ablehnung, die er je hören musste. Er nickte betrübt und machte Anstalten, zum nächsten Tisch zu gehen. »Ich danke Ihnen für Ihre diplomatische Ausdrucksweise.«

»Gerne. Und jetzt zieh Leine.«

\begin{center}
∞∞∞
\end{center}

»Die Leute mögen dich nicht, Fnbsrd«, rief Golandor ihm während des Essens noch einmal in Erinnerung. Es bestand wohl keine Aussicht darauf, das Thema zu wechseln.

»Das weiß ich doch«, seufzte Fnbsrd gelangweilt. »Man kommt sich fast vor, als fände diese Unterhaltung in einem Dorf statt, in dem jeder jeden kennt und noch mit Holz geheizt wird.«

Ein interessanter Punkt. »Vielleicht solltest du dorthin gehen«, schlug der Elf ernsthaft vor. »Vielleicht hat man dort mehr Verständnis.«

»Verständnis wofür?«, fragte Fnbsrd.

»Für jemanden, der mit dem Feuer spielt und seine Umgebung in Flammen setzt, aber keine Elementzauber zur Beherrschung der Katastrophe kennt.«

»Ich bin Künstler! Ich habe keine Muße und keinen Bedarf an Elementzaubern. Feuer, Wasser, Wind... Das sollen andere erledigen. Mein Metier ist die Schaffung, Veränderung und Reparatur von Leben, von Welten, von Weltsteinen.«

»Du bist ein Nichtsnutz. Zum Heiler hat es gerade so gereicht; mit brandgefährlichen Steinen ist kein Geld zu verdienen. Jeder, der ein Feuer entzünden möchte, kann das mit Grundschulzaubern in seinem Ofen tun. Du kannst nicht einmal das wirklich gut.«

»Die Grundschule verdirbt die Kinder«, echauffierte sich Fnbsrd. »Sie bereitet sie darauf vor, namenlos in ein System integriert, von diesem verschluckt und bedeutungslos zugrundegearbeitet zu werden. Am Ende steht ein Häufchen Asche, das als Dünger verwendet wird. Mich ekelt diese Gesellschaft an!«

»Oho, das beruht auf Gegenseitigkeit. Iss auf, ich habe nicht den ganzen Tag Zeit, dich zu beaufsichtigen.«

Fnbsrd explodierte innerlich. Nach außen zischte es. »Würde es dir etwas ausmachen, wenn ich mich dabei verschlucke?«

»Willst du eine ehrliche oder eine freundliche Antwort?«

\begin{center}
∞∞∞
\end{center}

Fnbsrd spielte seit Längerem mit dem Gedanken, Blautal den Rücken zu kehren und den Berg zu verlassen. Um die »Nabe«, die Bergspitze, herum gab es noch zwei weitere Dörfer: Grüntal, das Dorf eines blutrünstigen Elfenkönigs, und Rottal, das Dorf eines weisen Drachenherrschers. Blautal wurde von einem scheuen Menschen regiert, der in seiner Angst vor Fremden unzählbare Seelen auf dem Gewissen hatte. Sprechenderweise hieß der Elf »Thürstön«, der Drache »instein« und der Mensch »Xenophob«. Die drei Könige der Talstädte kamen dem, was manche sich unter Göttern vorstellen, sehr nah; ihr Kontakt zur Außenwelt war seit Langem durch Magie und fehlenden physischen Kontakt geprägt. Ihre Unsterblichkeit bezogen sie durch Nekromantie aus dem Leben ihrer Untertanen. Ein paar Jahre weniger Lebenszeit für jeden Einzelnen fielen kaum ins Gewicht; den Königen bescherte es das Überdauern der Jahrtausende. Früher in Kriegen, heute in Angst. Kalt und abstoßend war das alles.

Um den Berg herum führten niedere, schlecht entwickelte Dörfer als Stellvertreter den Krieg der Herren untereinander. Wer in der Gunst des Grüntals kämpfte, erhielt vom Elfenkönig und seinen Untertanen Waffen zur Eroberung der unendlich großen Welt. Wer im Dienst des Drachen ein Dorf regierte, konnte mit strategischer Unterstützung rechnen, die ihm einen Vorteil gegenüber den wenigen gunstlosen Dörfern verschaffte, die in ihrer »Neutralität« stets am Abgrund standen. Der menschliche König Xenophob war ein Defensivkünstler, ein Mauernbauer, ein Burgarchitekt. Große Stadtmauern, die so nie von Hand einzelner Lebewesen erschaffen werden konnten, umgaben die von ihm unterstützten Kleinstädte.

Man konnte bei dieser Aufteilung als Fremder leicht auf den Gedanken kommen, in Blautal lebten hauptsächlich Menschen, im Grüntal hauptsächlich Elfen und im Rottal fast nur Drachen. Dem war nicht so. Und wer bei Magie und Drachen glaubte, diese seien im Umgang mit Feuer begabter als Elfen, der irrte ebenfalls gewaltig. Jedes Kind konnte zaubern, jeder Jugendliche war auf irgendeine Magieform spezialisiert, jeder Erwachsene war ein Erzmagier seiner Schule. Talent und Fleiß beim Üben der Zauber, die Wahl der Sprüche und der Spezialisierung waren niemandem durch seine Abstammung in die Wiege gelegt. Es gab talentiert wasserzaubernde Drachen ebenso wie Windmagier unter den Menschen. Trolle, Zwerge, viele mystische Wesen bevölkerten die magiegetränkte Welt, und keines von ihnen war intrinsisch zu einer Form der Magie gezwungen.

Fnbsrd hatte die Erschaffung von Weltsteinen für sich entdeckt. Schwarz wie seine Kleidung, aber in unsichtbar tiefem Rot glühend wie das Innere von Vulkanen, nur kompakter und gebändigter in den Händen erfahrener Künstler. Eine brotlose und gefährliche Kunst, denn Holz und die Steine vertrugen sich nicht. Mit modernen Portalen konnte man in diesen Steinen versinken, während draußen fast keine Zeit verging. Ein nahezu beliebig langer Ausflug in eine Gedankenwelt mit wirklich beliebig erschaffbaren Regeln und Charakteren. Der ideale Schlüssel zu unendlicher Weisheit, wenn man diese Steine zu nutzen wusste. Fnbsrd war gebildeter als die meisten Lehrer, überheblicher als ein König und ärmer als manche Bettler. Ein paar Heilzauber lagen in seinem Repertoire, vernachlässigt zugunsten der Kunst. Überlebensfähig war er durch das Mitleid seiner Mitmenschen, doch das hielt sich in Grenzen, wenn er zum wiederholten Male ganze Landstriche dem Feuer zum Opfer fallen ließ, nur weil er gedankenverloren einen Stein falsch abgelegt hatte.

Und so kam es, wie es kommen musste: Fnbsrd kehrte Blautal den Rücken und trat den Weg in die Tiefe am Fuß des Berges an. Was als »Talstadt« bekannt war, lag tausende Häuser hoch über dem Pöbel. Eine abgehobene Gesellschaft, die in ihrem Fortschritt versank und jede Sympathie verspielte.

»Auf Nimmerwiedersehen«, war der Abschiedsgruß, zu dem sich wenigstens ein Elf herabließ. Es war Golandor, derjenige, der Fnbsrd gegenüber am wenigsten negativ eingestellt war. Sozusagen ein bester Freund. Alle anderen schickten stumm ein paar Flüche hinterher oder dankten dem König dafür, dass er sie von der Last erlöst hatte.

»Auf Nimmerwiedersehen. Ich werde euch nicht vermissen.«

Golandor musste das letzte Wort haben. »Wir dich auch ganz bestimmt nicht.«

\begin{center}
∞∞∞
\end{center}

Die Dorfmauern von Last Hope waren eine Eigenkonstruktion der darin lebenden Bewohner. Lehm, billige Ziegel, unregelmäßige Zinnen, aber immerhin bis zu drei Meter hoch über den Boden ragend. Einen Burggraben gab es nicht; das Gebilde als »Burg« zu bezeichnen, war eine Beleidigung gegenüber den von Xenophob magisch unterstützten Architekturen in anderen Dörfern. Doch von Xenophob und Konsorten wollte hier niemand etwas wissen; Könige waren unwillkommenes Gesindel vom Berg.

\iathought{Beste Voraussetzungen für einen Besuch} sah Fnbsrd darin. Er klopfte von außen an das riesige hölzerne Doppeltor und ersuchte um Einlass.

»Wer sind Sie und was wollen Sie?«, schallte es ihm von oben entgegen.

Fnbsrd reckte seinen Kopf in die Höhe und rief zurück: »Ich bin die Dunkelheit und suche nach Licht.«

Das war ungewöhnlich genug, um eine Rückfrage bei einem Dienstälteren hervorzurufen. Mangels Telefon und Privatsphärebedürfnis schallte diese einfach über den Platz. Aus der Ferne hörte man eine eher murrende Antwort, dann kamen Schritte näher. Ein roter Drache, etwa zweieinhalb Meter groß %TODO: In Ellen umwandeln. Umwandlungsfaktor definieren (Bonusmaterial).
und dennoch weit vom Torbogen entfernt, drückte mit seinen Klauen das Holz beiseite. Fnbsrd ahnte, dass Magie hinter der vorgeblichen Kraftdemonstration steckte, und zeigte sich entsprechend unbeeindruckt. »Moin.«

»Wer sind Sie?«, fragte der Drache mit tiefer, dröhnender Stimme. Aus seinen Nüstern züngelte das innere Feuer.

»Mein Name ist Fnbsrd.« Er sprach jeden Buchstaben einzeln aus, weil die Konsonanten sich absichtlich nur buchstabieren oder durch unanständige Vokale ergänzen ließen. Die Namenswahl war eine der vielen Gemeinheiten, mit denen man im Blautal sogar von seinen Erzeugern gestraft wurde.

Der Drache schnaubte vorsichtig. »Ein ungewöhnlicher Name. Und keine Antwort auf meine Frage.«

Fnbsrd blinzelte.

»Woher kommen Sie?«

»Blautal.«

Feuer!

»Brauchen Sie ein Taschentuch?«

Fnbsrd stand in einer Flammenwolke; um ihn herum verbrannte die spärliche Vegetation zu schwarzer Asche. Der Drache blinzelte; das Blatt hatte sich gewendet.

»Hören Sie, ich suche Zuflucht vor den Idioten meiner Ex-Stadt. Mir sind Ihre Mauern bereits von außen sympathisch; wenn das Dorfinnere nur die Hälfte dieses Willkommens verstrahlt, möchte ich gerne hier bleiben.«

Der Junge war in Ordnung. Ein Wunder, dass er noch lebte. »Ich bitte um Verzeihung… Fnbsrd?«

Fnbsrd nickte.

»Mein Name ist curie. Ich bin Flammenmagierin und bilde mich derzeit in Erdkunde fort.«

\iathought{Eine sinnvolle Umschulung.} »Sehr erfreut.«

»Die Freude ist ganz meinerseits. Bitte treten Sie ein.« Der Drache machte eine einladende Geste und trat ein paar Schritte zur Seite. Neben ihm hätten zwar mindestens zwei weitere Drachen in das Tor gepasst, aber die Höflichkeit gebot ein sanftes Ausweichen.

Das ließ sich Fnbsrd nicht zweimal sagen. Er verließ die Vegetationsreste und betrat eine abgeschlossene Welt des Friedens.


\chapter{Last Hope}

Fnbsrd war es aus seiner Heimat gewohnt, dass körperliche Arbeit vollständig durch Magie automatisiert wurde. Niemand pflügte einen Acker, niemand trug Wasser aus Brunnen in Häuser.

curie bot dem Gast eine kleine Dorfführung. »Wir haben einem Teil dessen, was ihr dort oben als ›Fortschritt‹ betrachtet, vor die Burgmauern verbannt.« Der Drache winkte einem älteren Menschenherrn zu, der bunte Früchte in einer Holzkiste transportierte. Neben ihm trug ein Esel geduldig zwei Säcke Mehl über dem Sattel.

»Was spricht gegen ein wenig Luftmagie?«, wunderte sich der Junge.

»Dekadenz« war das Stichwort. »Womit verbringt ihr eure Lebenszeit?«

»Mit Arbeit.« D'oh.

Der Drache war nicht weniger amüsiert als sein Nebenmann. »Dienstleistungen.«

»Ein starker Tertiärsektor ist charakteristisch für Fortschritt«, brabbelte Fnbsrd. Dabei war er eigentlich eher unter den Systemkritikern zu verordnen. Hier jedoch setzte sich der Impuls durch, seine Herkunft gegen Unverständnis zu verteidigen. »Was der Herr dort tut, ist doch eines Intelligenzwesens nicht würdig.«

»Möchten Sie einen Apfel kaufen?«, fragte der nicht von Schwerhörigkeit betroffene Mann freundlich. Er blickte Fnbsrd tief durch die Sonnenbrille und versank in ewiger Finsternis. »Whoa. Sie haben schwarze Augen.«

»Ja, und gerne.« Der Besucher grinste den Verkäufer frech an. »Fünf Xennys?«

Ein Spiegel? Das Grinsen kam jedenfalls zurück. »Drei Holzhelax.«

»Deine Königswährung ist hier nichts wert, Fremder«, erläuterte curie das nun Offensichtliche noch einmal zur Sicherheit. »Old Jonny macht dir einen Freundschaftspreis, aber du bist mittellos.«

Das war unangenehm, aber erwartbar gewesen. »Xennys bestehen aus Kupfer«, wagte Fnbsrd einen Wechselversuch. »Bedeutet euch das etwas?«

»Mäßig.« Der Drache zog umständlich einen großen Geldsack zwischen seinen Rückenschuppen hervor. Dann ging er zu einem öffentlichen Tisch und verteilte einige Münzen darauf. Viel Holz war dabei, mäßig Kupfer, wenig Silber und eine Goldmünze. »Bist du gut in Mathematik?«

\iathought{Sehr.} »Ja, doch schon.« \iathought{Falls du mich über den Tisch ziehen möchtest, rechne ich dich zu Boden.}

curie freute sich. »Ein Goldhelax sind zehn Silberhelax. Ein Silberhelax sind zehn Kupferhelax. Eigentlich ganz simpel. Holzhelax fallen aus der Reihe; hundert von ihnen ersetzen einen Kupferhelax.«

Der »Freundschaftspreis« schien tatsächlich einer zu sein, selbst bei knausriger Umrechnung. Unter »Xennys« verstand man im Blautal die kleinste Währungseinheit, während die aus Messing gefertigten »Xenos« jeweils hundert Xennys darstellten. Große Geldtransfers fanden nur über Banken statt; zur Not behalf man sich mit Papierxenos, die magische Hologramme trugen und durch das Bankenwesen zunehmend aus dem Umlauf verschwanden. Nichts dergleichen gab es hier im Dorf?

»Kannst du mit dem Begriff ›Buchgeld‹ etwas anfangen?«, fragte Fnbsrd vorsichtig.

»Du musst noch viel lernen«, spottete die angehende Geologin. »Erinnere dich an deine Frage, wenn du eine Schubkarre voll Helax vor dir herschiebst, und höre dann mein Nein.«

Niemand schob Schubkarren durch das Dorf. So weit man blickte, war das ein Fantasiebild.

»Du suchst nach Schubkarren«, stellte der Obstverkäufer fest. »Mein Name ist übrigens Jonathan.«

»Fnbsrd.« Die beiden reichten sich die Hände. »Richtig. Ich sehe aber keine.«

»Jetzt hast du den Witz verstanden.«

Der Drache fand das offensichtlich lustig; Fnbsrd verkrampfte leicht. Das fing ja gut an. Kleine Flammen schossen im Lachtakt aus der Feuernase.

»Wie komme ich nun an Helax?«, fragte der Ankömmling leicht genervt.

»Die drei Holzhelax für den Apfel kannst du geschenkt haben; ich entlaste dich im Gegenzug von deiner wertlosen Währung.«

Bei Fnbsrd schrillten die Alarmglocken. »Kupfer ist wertlos?«

»Dort hinten gibt es einen Schmied, der kauft es dir für ein paar Holzhelax ab«, gab curie zu. »Versprich dir aber keinen Reichtum davon. Wie du siehst, prägen auch wir damit nicht unsere wertvollsten Münzen.«

»Hm.« Nachdenklich blickte sich Fnbsrd entlang der Stadtmauern um. »Wie groß ist das Dorf eigentlich?«

»Knapp vier Quadratkilometer.« Das war keine Schätzung. curie kannte die Maße auswendig. %TODO: In Ellen umwandeln, Umwandlungsfaktoren definieren und im Bonusteil dokumentieren

Die Frage nach der Einwohnerzahl stand vorhersehbar im Raum; Jonathan lachte kratzend. »Für tausend Bewohner sehr geräumig. Wir sind von Expansionsgedanken weit entfernt.«

»Ihr habt die Mauern von Hand errichtet, nicht wahr?«

Nun lachten curie und Jonathan recht laut. »Sieht man das nicht?«

»Es war wohl eine rhetorische Frage«, murrte Fnbsrd. »Ihr habt euren Spaß, ich habe nichts zu essen.«

Das ließ Old Jonny nicht auf sich sitzen. »Hör mal, Kleiner. Iss dich satt. Die Bäume sind Gemeinschaftseigentum.« Er wies in Richtung eines großen Parks. »Du musst sie nicht von mir kaufen.«

Nun verstand der wirtschaftlich nicht ungebildete Junge die Welt nicht mehr. »Aber wieso werden sie überhaupt gekauft?«

»Du musst noch viel lernen«, wiederholte curie ihr Mantra. »Wenn du magst, bezahle ich dich fürs Umgraben der Erde. Fünf Holzhelax pro Stunde für eine stumpfe, sinnlose Tätigkeit, mit der du deinen Tag verbringen kannst.«

»Ah«, machte Fnbsrd. »Es gibt bei euch keine Arbeitslosigkeit.«

»Es gibt bei uns keinen Arbeitsmangel, du Held. Wer arbeiten möchte, findet Arbeit. Wer nicht arbeitet, langweilt sich zu Tode.«

Das war hoffentlich nicht wörtlich gemeint.

»Das meine ich wörtlich.«

\begin{center}
∞∞∞
\end{center}

Was zunächst wie ein Wechsel vom Regen in die Traufe wirkte, ergab bei längerer Überlegung Sinn. Zumindest aus Sicht der Dorfbewohner war die Situation begrüßenswert; es gab keinen Bedarf, das Gleichgewicht zu verändern. Was Fnbsrd als »Fortschritt« kannte, war hier verpönt – aus Erfahrung oder Beobachtung, jedenfalls einigermaßen empirisch.

Aus dem tiefen Bedürfnis heraus, nicht mittellos neue Kontakte zu knüpfen, ließ sich Fnbsrd für eine Stunde sinnloses Graben bezahlen. Zudem verkaufte er sein Kupfer und stand am Ende mit dreizehn Holzmünzen in der Hand vor der Tür einer Gaststätte. »Zur ewigen Ruhe« hieß es auf einem schwarzen Schild mit weißer Beschriftung, die aussah, als wäre sie mit wasserlöslicher Farbe im Regen geschrieben worden, bevor jemand das Schild aus Mitleid lackiert hatte. Nun war es offenbar wasserfest.

»Öffnet mir, lasst mich hinein«, bat Fnbsrd. Er schlug nun zum wiederholten Mal mit einem rostigen Klopfer gegen das Eisen an der Tür. Es war zwar lange noch nicht Nacht, aber der Regen begünstigte den Wunsch nach einem eigenen Zimmer.

»Meine Tür versperrt ein Eisenschloss«, klang es dann von drinnen. War das die Stimme eines Elfen?

Er hatte jedenfalls ein Talent dazu, Unnötiges zu erwähnen. »Das ist mir nicht entgangen!«

»Ich habe keinen Schlüssel dafür.« Schließlich öffnete er die Tür doch. Das war so unlogisch wie das Lied »Diese kalte Nacht« von Faun. Einfach nicht hinterfragen. Jedenfalls stand dort tatsächlich ein Elf, vielleicht 1,80 Meter groß, blond, klischeehaft grünbraun gekleidet wie der Wald.

»Äh, vermieten Sie Zimmer?«, fragte der Junge mit dem Talent für überflüssige Fragen. Da hatten sich zwei gefunden.

»Hin und wieder.« Der Elf stellte sich als »Öbfüs« vor, nahm widerwillig die münzgewordene Armut entgegen und mehrte sie dadurch wie ein Profi. »Du kannst einen halben Mond lang hier bleiben.«

Fnbsrd wollte sich bedanken, doch er wurde unterbrochen.

»Falls es wirklich sein muss.«

\iathought{Sehr freundlich.} »Es muss. Es muss.«

»Dann hereinspaziert und willkommen in der Hölle.«

Der Elf hatte einen sehr eigenartigen Humor, aber Fnbsrd kam gut damit zurecht. Von seiner Dunkelheit konnte sich der Wirt ein paar Scheiben abschneiden. So hatten sich wirklich zwei Gleichgesinnte gefunden.

»Wir haben keinen Teer mehr im Kühlschrank, aber ich kann dir ein Glas Wasser anbieten.«

Reizend. »Ich nehme zwei.«

Das war eine gute Antwort, musste Öbfüs zugeben. »Hör mal, deine Lebenseinstellung gefällt mir.«

»Schade.« Bäm.

»Guck bei Gelegenheit mal in der Kneipe an der östlichen Burgmauer vorbei. Ich glaube, dort hängt gerade etwas für dich am schwarzen Brett aus.« War das Absicht? Der Elf grinste. Ja, das war genau so beabsichtigt.


\chapter{To Hell and Back}

»Zur Hölle und wieder zurück.« Moment. Nein, da stand nichts von einer Rückreise. »Zur Hölle.«

Das schwarze Brett bot neben einem Arbeitsüberangebot im Dorf auch Gesuche von außen. Eines davon hatte es in sich.

»›Lebensmüder Abenteurer ohne Ortsbindung und emotionale Abhängigkeiten zur Vervollständigung eines vierköpfigen Todeskommandos gesucht.‹ Nehm ich.«

Die Entscheidung war bereits beim Lesen gefallen. Hier bot sich Abwechslung, Umbruch und Reichtum zugleich. Reichtum an Erfahrung, da Fnbsrd die Wahrscheinlichkeit für eine Rückkehr mit Goldmünzen naiv überschätzte.

Viele der anderen Zettel am Brett zeichneten sich dadurch aus, dass die abtrennbaren Kontaktinformationen mindestens angebrochen, teils aufgebraucht waren. In unerhörter Eigenmacht riss Fnbsrd nicht nur den ersten Kontaktzettel von dem Höllenangebot ab, sondern entfernte zudem alle kontaktlosen Angebote vom Brett. Die Aushänge hatten ihren Zweck erfüllt und waren nun aufgebraucht, das ergab also durchaus Sinn. Es war nur ungewöhnlich, dass ein Gast so offen als Verwalter auftrat.

»Hat dir das jemand erlaubt?«, fragte daher noch während des Prozesses eine tiefe Stimme von hinten.

Jeder andere wäre herumgefahren und hätte sich vorgestellt. Fnbsrd riss weitere Zettel ab und murmelte seinen Namen, ohne den Blick vom Brett zu nehmen. Schließlich wurde es dem Troll in seinem Rücken zu bunt und eine steinschwere inflexible Pranke legte sich auf die rechte Schulter des Menschenjungen.

Fnbsrd beschloss, seine Arbeit für eine kurze Nachricht zu unterbrechen. Er wandte seinen Kopf gefühlt um 180 Grad, spürte ein Nachlassen des Schulterdrucks und nutzte dieses, um ruckartig voll herumzufahren. »Besteht noch Bedarf an den vollständig abgerissenen Angeboten?«

»Äh, nee.« Einige Schaulustige unterbrachen ihr Mahl und blickten verstohlen herüber. Die Geräuschkulisse erstickte.

»Dann führ nicht so ein Drama auf, wenn jemand sie entfernt.«

»Kennen wir uns?!«, fragte der Troll erbost. Es handelte sich, wie spätestens jetzt auch Fnbsrd klar geworden sein musste, um den Herrn des Hauses.

»Keine Sorge, du wirst mich noch kennenlernen.« Zwei weitere Zettel waren schneller abgerissen, als der Troll gucken konnte. Fnbsrd drückte ihm das gesammelte Altpapier in die Hand. »Bitte entsorgen. Schönen Tag noch.«

Mit diesen Worten verließ Fnbsrd die Ostkneipe.

\begin{center}
∞∞∞
\end{center}

Öbfüs bekam sich vor Lachen kaum noch ein. Gemeinsam mit Fnbsrd saß er an einem kleinen Esstisch; es wurde Vollkornbrot in Tomatensauce mit Basilikum geboten. »Das geschieht Rambo ganz recht. Der nimmt sich viel zu wichtig.«

»Seine Worte haben Gewicht«, räumte Fnbsrd schmunzelnd ein.

»Eine Pfeife ist das. Weil er sie auf der richtigen Seite gebaut hat, ist seine Kneipe ein vielbesuchter Ort. Das liegt aber nicht an toller Planung, sondern an einem nachträglich eingebauten Dorftor. Und dorthin führt inzwischen ein kaum noch als Schleichweg bezeichenbarer Pfad von anderen Dörfern.«

»Sind die anderen Dörfer auch so konservativ gegenüber Hochbildung und automatisierter Arbeit?«

»He, wir haben gute Bibliotheken und eine gute Schule.«

»Schön«, gab Fnbsrd zu. »Und wie lautet die Antwort auf meine Frage?«

»Du würdest jede Umschreibung mit einem simplen ›Ja‹ zusammenfassen, also bekommst du es gerne direkt.«

»Du klingst wie ein dämlicher Städter, der mich im Blautal angefeindet hat.«

»Bist du hier, weil du Heimweh hast?«, fragte Öbfüs spitz. Das saß. Sanfter fuhr er fort: »Wir sind stolz darauf, nicht so zu sein wie die Menschen vom Berg. Daher fühlst du dich hier ja auch so wohl. Und curie hätte dich kaum ins Dorf gelassen, wenn du zur Stadt gepasst hättest.«

Fnbsrd gab zu, dass diese Gedankengänge auch für ihn sehr nachvollziehbar waren. »Wir sind da auf einer Wellenlänge.«

»Davon verstehe ich nichts.«

»Was ich sagen wollte: Für uns spielt dasselbe Schlagzeug.«

»Ah. Das ist eine schöne Metapher.«

»Gibt es Berufsmusiker bei euch?«

»Nein, auch wenn man bei uns durchaus davon leben könnte. Alle wollen etwas tun, das über geistiges Schaffen hinausgeht.«

Fnbsrd, für den fast unverrückbar geistiges Schaffen die höchste Entwicklung intelligenter Lebewesen darstellte, hielt dies für eine unlogische Aussage.

»Wir möchten am Ende des Tages spüren, dass wir gearbeitet haben. Viele von uns zumindest.«

»Und das tust du, Gastwirt?«

Öbfüs legte einen Ellenbogen auf den Tisch. »Es gibt in diesem Haus nur ein Bett und der Verlierer schläft auf dem Boden.«

Fnbsrd lachte ihn auf der Stelle aus. Armdrücken? »Hier hast du deine Gläser zurück. Vielleicht kannst du sie als Kopfkissen verwenden.«

Dem Jungen war mit Witz und Kraft nicht beizukommen. Natürlich gab es Gästebetten. Eher mit Freude über einen so schlagfertigen Gast als mit Frustration räumte Öbfüs die Gläser ins Spülbecken. Als der Abwasch erledigt war, spürte Öbfüs, dass er gearbeitet hatte, und legte sich im Erdgeschoss zu Bett. Fnbsrd schlief zu diesem Zeitpunkt bereits tief und fest unter einer Dachschräge auf seiner Bettdecke.

\begin{center}
∞∞∞
\end{center}





































































»Um Himmels willen!«

In jeder Facette seiner Ambivalenz der richtige Ausspruch.












\part{Bonusmaterial}

\chapter {Postskriptum}

(noch nicht geschrieben)


\chapter{Titelmelodie}

\begin{figure}[p]
    \includegraphics[width=\textwidth, page=1]{z-include-fnbsrdtheme.pdf}
\end{figure}


\chapter{Musikliste}

\textbf{Für Filmproduzenten, Träumer und Multitasking-Genies.}

\begin{itemize}
    \item Falls Du ernsthaft einen Film zu diesem Buch drehen möchtest.
    \item Falls Du das gesamte Buch bereits ausgelesen hast und die genannten Lieder vielleicht noch nicht kennst. Höre die Lieder und stelle Dir dabei die Szenen vor. Wenn es schon keinen Fnbsrd-Film gibt, kannst Du wenigstens einen Film in deinem Kopf laufen lassen.
    \item Falls Du beim Lesen Musik hören möchtest, die zur aktuellen Szene passt.
\end{itemize}

Diese Liste wurde von Tobias Frei zusammengestellt und impliziert keinerlei Unterstützung oder Befürwortung durch die Komponisten der Lieder. Eines Tages wird jedes dieser Lieder in die Gemeinfreiheit übergehen; der genaue Zeitpunkt hängt von verschiedenen Gesetzen ab.

\begin{enumerate}
    \item Titelmelodie:\\ »Fnbsrd«~– Tobias »ToBeFree« Frei
    \item \textbf{Teil 1: Fnbsrd.}\\ »Träumen«~– Solitary Experiments featuring In Strict Confidence
    \item Unwillkommenheit:\\ »Spectres from the Black Moss«~– Ashbury Heights
    \item Auf Nimmerwiedersehen:\\ »Oh Johnny«~– Jan Delay
    \item Vor den Mauern der letzten Hoffnung:\\ »Storm«~– Assemblage 23
    \item Willkommen in Last Hope:\\ »Songs of Love and Death«~– Beyond The Black
    \item Wirtschaftslektion am Tisch:\\ »Crash and Burn«~– Solitary Experiments
    \item Fnbsrd vor der Tür:\\ »Diese kalte Nacht«~– Faun
    \item Am schwarzen Brett:\\ »To Hell and Back«~– Sabaton (Acoustic Cover by Beyond the Black)
    \item Erster Kampf:\\ »Infinite«~– Assemblage 23
\end{enumerate}

Musik, die leider nicht mehr in die Musikliste gepasst hat:

\begin{itemize}
    \item »Lorem«~– Ipsum
\end{itemize}


\chapter{Bildquellen}

Alle verwendeten Bilder sind gemeinfrei. Die Verwendung der Bilder in diesem Roman impliziert keinerlei Unterstützung oder Befürwortung durch ihre Schöpfer.

\begin{itemize}
    \item Derzeit werden keine Bilder verwendet.
\end{itemize}


\chapter{Lizenz des Buchinhalts}

\textbf{Fnbsrd © by\\ Tobias Frei, fnbsrd.de}

Dies ist eine offizielle Ausgabe des Romans Fnbsrd, herausgegeben von Tobias Frei. Veränderte Versionen und unautorisierte Nachdrucke müssen deutlich als solche erkennbar sein. Auch das Impressum muss angepasst werden, wenn das Dokument verändert wird.

Falls du die Rechte in dieser Lizenz nutzen möchtest, musst du sie vollständig gelesen und verstanden haben. Es genügt nicht, nur eine Zusammenfassung zu lesen. Aus diesem Grund wird in diesem Buch keine Zusammenfassung angeboten.

This novel is licensed under a Creative Commons Attribution-ShareAlike 4.0 International License.

You should have received a copy of the license along with this work. If not, see\\
https://creativecommons.org/licenses/by-sa/4.0/legalcode

You are required to actually read and understand the full text of the license, not a summary.

\begin{center}
    \large{\textbf{Creative Commons Attribution-ShareAlike 4.0 International Public License}}
\end{center}

By exercising the Licensed Rights (defined below), You accept and agree to be bound by the terms and conditions of this Creative Commons Attribution-ShareAlike 4.0 International Public License ("Public License"). To the extent this Public License may be interpreted as a contract, You are granted the Licensed Rights in consideration of Your acceptance of these terms and conditions, and the Licensor grants You such rights in consideration of benefits the Licensor receives from making the Licensed Material available under these terms and conditions.

\begin{center}
    \textbf{Section 1 -- Definitions.}
\end{center}

\begin{itemize}
    \item[a.] \textbf{Adapted Material} means material subject to Copyright and Similar Rights that is derived from or based upon the Licensed Material and in which the Licensed Material is translated, altered, arranged, transformed, or otherwise modified in a manner requiring permission under the Copyright and Similar Rights held by the Licensor. For purposes of this Public License, where the Licensed Material is a musical work, performance, or sound recording, Adapted Material is always produced where the Licensed Material is synched in timed relation with a moving image.
    \item[b.] \textbf{Adapter's License} means the license You apply to Your Copyright and Similar Rights in Your contributions to Adapted Material in accordance with the terms and conditions of this Public License.
    \item[c.] \textbf{BY-SA Compatible License} means a license listed at creativecommons.org/compatiblelicenses, approved by Creative Commons as essentially the equivalent of this Public License.
    \item[d.] \textbf{Copyright and Similar Rights} means copyright and/or similar rights closely related to copyright including, without limitation, performance, broadcast, sound recording, and Sui Generis Database Rights, without regard to how the rights are labeled or categorized. For purposes of this Public License, the rights specified in Section 2(b)(1)-(2) are not Copyright and Similar Rights.
    \item[e.] \textbf{Effective Technological Measures} means those measures that, in the absence of proper authority, may not be circumvented under laws fulfilling obligations under Article 11 of the WIPO Copyright Treaty adopted on December 20, 1996, and/or similar international agreements.
    \item[f.] \textbf{Exceptions and Limitations} means fair use, fair dealing, and/or any other exception or limitation to Copyright and Similar Rights that applies to Your use of the Licensed Material.
    \item[g.] \textbf{License Elements} means the license attributes listed in the name of a Creative Commons Public License. The License Elements of this Public License are Attribution and ShareAlike.
    \item[h.] \textbf{Licensed Material} means the artistic or literary work, database, or other material to which the Licensor applied this Public License.
    \item[i.] \textbf{Licensed Rights} means the rights granted to You subject to the terms and conditions of this Public License, which are limited to all Copyright and Similar Rights that apply to Your use of the Licensed Material and that the Licensor has authority to license.
    \item[j.] \textbf{Licensor} means the individual(s) or entity(ies) granting rights under this Public License.
    \item[k.] \textbf{Share} means to provide material to the public by any means or process that requires permission under the Licensed Rights, such as reproduction, public display, public performance, distribution, dissemination, communication, or importation, and to make material available to the public including in ways that members of the public may access the material from a place and at a time individually chosen by them.
    \item[l.] \textbf{Sui Generis Database Rights} means rights other than copyright resulting from Directive 96/9/EC of the European Parliament and of the Council of 11 March 1996 on the legal protection of databases, as amended and/or succeeded, as well as other essentially equivalent rights anywhere in the world.
    \item[m.] \textbf{You} means the individual or entity exercising the Licensed Rights under this Public License. Your has a corresponding meaning.
\end{itemize}

\begin{center}
    \textbf{Section 2 -- Scope.}
\end{center}

\begin{itemize}
    \item[a.] \textbf{License grant.}
    \begin{itemize}
        \item[1.] Subject to the terms and conditions of this Public License, the Licensor hereby grants You a worldwide, royalty-free, non-sublicensable, non-exclusive, irrevocable license to exercise the Licensed Rights in the Licensed Material to:
        \begin{itemize}
            \item[A.] reproduce and Share the Licensed Material, in whole or in part; and
            \item[B.] produce, reproduce, and Share Adapted Material.
        \end{itemize}
        \item[2.] \underline{Exceptions and Limitations}. For the avoidance of doubt, where Exceptions and Limitations apply to Your use, this Public License does not apply, and You do not need to comply with its terms and conditions.
        \item[3.] \underline{Term}. The term of this Public License is specified in Section 6(a).
        \item[4.] \underline{Media and formats; technical modifications allowed}. The Licensor authorizes You to exercise the Licensed Rights in all media and formats whether now known or hereafter created, and to make technical modifications necessary to do so. The Licensor waives and/or agrees not to assert any right or authority to forbid You from making technical modifications necessary to exercise the Licensed Rights, including technical modifications necessary to circumvent Effective Technological Measures. For purposes of this Public License, simply making modifications authorized by this Section 2(a)(4) never produces Adapted Material.
        \item[5.] \underline{Downstream recipients}.
        \begin{itshape}\begin{itemize}
            \item[A.] \underline{Offer from the Licensor -- Licensed Material}. Every recipient of the Licensed Material automatically receives an offer from the Licensor to exercise the Licensed Rights under the terms and conditions of this Public License.
            \item[B.] \underline{Additional offer from the Licensor -- Adapted Material}. Every recipient of Adapted Material from You automatically receives an offer from the Licensor to exercise the Licensed Rights in the Adapted Material under the conditions of the Adapter's License You apply.
            \item[C.] \underline{No downstream restrictions}. You may not offer or impose any additional or different terms or conditions on, or apply any Effective Technological Measures to, the Licensed Material if doing so restricts exercise of the Licensed Rights by any recipient of the Licensed Material.
        \end{itemize}\end{itshape}
        \item[6.] \underline{No endorsement}. Nothing in this Public License constitutes or may be construed as permission to assert or imply that You are, or that Your use of the Licensed Material is, connected with, or sponsored, endorsed, or granted official status by, the Licensor or others designated to receive attribution as provided in Section 3(a)(1)(A)(i).
    \end{itemize}
    \item[b.] \textbf{Other rights.}
    \begin{itemize}
        \item[1.] Moral rights, such as the right of integrity, are not licensed under this Public License, nor are publicity,
          privacy, and/or other similar personality rights; however, to the extent possible, the Licensor waives and/or agrees not to assert any such rights held by the Licensor to the limited extent necessary to allow You to exercise the Licensed Rights, but not otherwise.
        \item[2.] Patent and trademark rights are not licensed under this Public License.
        \item[3.] To the extent possible, the Licensor waives any right to collect royalties from You for the exercise of the Licensed Rights, whether directly or through a collecting society under any voluntary or waivable statutory or compulsory licensing scheme. In all other cases the Licensor expressly reserves any right to collect such royalties.
    \end{itemize}
\end{itemize}

\begin{center}
    \textbf{Section 3 -- License Conditions.}
\end{center}

Your exercise of the Licensed Rights is expressly made subject to the following conditions.

\begin{itemize}
    \item[a.] \textbf{Attribution.}
    \begin{itemize}
        \item[1.] If You Share the Licensed Material (including in modified form), You must:
        \begin{itemize}
            \item[A.] retain the following if it is supplied by the Licensor with the Licensed Material:
            \begin{itemize}
                \item[i.] identification of the creator(s) of the Licensed Material and any others designated to receive attribution, in any reasonable manner requested by the Licensor (including by pseudonym if designated);
                \item[ii.] a copyright notice;
                \item[iii.] a notice that refers to this Public License;
                \item[iv.] a notice that refers to the disclaimer of warranties;
                \item[v.] a URI or hyperlink to the Licensed Material to the extent reasonably practicable;
            \end{itemize}
            \item[B.] indicate if You modified the Licensed Material and retain an indication of any previous modifications; and
            \item[C.] indicate the Licensed Material is licensed under this Public License, and include the text of, or the URI or hyperlink to, this Public License.
        \end{itemize}
        \item[2.] You may satisfy the conditions in Section 3(a)(1) in any reasonable manner based on the medium, means, and context in which You Share the Licensed Material. For example, it may be reasonable to satisfy the conditions by providing a URI or hyperlink to a resource that includes the required information.
        \item[3.] If requested by the Licensor, You must remove any of the information required by Section 3(a)(1)(A) to the extent reasonably practicable.
    \end{itemize}
    \item[b.] \textbf{ShareAlike.}

     In addition to the conditions in Section 3(a), if You Share Adapted Material You produce, the following conditions also apply.

    \begin{itemize}
        \item[1.] The Adapter's License You apply must be a Creative Commons license with the same License Elements, this version or later, or a BY-SA Compatible License.
        \item[2.] You must include the text of, or the URI or hyperlink to, the Adapter's License You apply. You may satisfy this condition in any reasonable manner based on the medium, means, and context in which You Share Adapted Material.
        \item[3.] You may not offer or impose any additional or different terms or conditions on, or apply any Effective Technological Measures to, Adapted Material that restrict exercise of the rights granted under the Adapter's License You apply.
    \end{itemize}
\end{itemize}

\begin{center}
    \textbf{Section 4 -- Sui Generis Database Rights.}
\end{center}

Where the Licensed Rights include Sui Generis Database Rights that apply to Your use of the Licensed Material:

\begin{itemize}
    \item[a.] for the avoidance of doubt, Section 2(a)(1) grants You the right to extract, reuse, reproduce, and Share all or a substantial portion of the contents of the database;
    \item[b.] if You include all or a substantial portion of the database contents in a database in which You have Sui Generis Database Rights, then the database in which You have Sui Generis Database Rights (but not its individual contents) is Adapted Material, including for purposes of Section 3(b); and
    \item[c.] You must comply with the conditions in Section 3(a) if You Share all or a substantial portion of the contents of the database.
\end{itemize}

For the avoidance of doubt, this Section 4 supplements and does not replace Your obligations under this Public License where the Licensed Rights include other Copyright and Similar Rights.

\begin{center}
    \textbf{Section 5 -- Disclaimer of Warranties and Limitation of Liability.}
\end{center}

\begin{itemize}
    \item[\textbf{a.}] \textbf{Unless otherwise separately undertaken by the Licensor, to the extent possible, the Licensor offers the Licensed Material as-is and as-available, and makes no representations or warranties of any kind concerning the Licensed Material, whether express, implied, statutory, or other. This includes, without limitation, warranties of title, merchantability, fitness for a particular purpose, non-infringement, absence of latent or other defects, accuracy, or the presence or absence of errors, whether or not known or discoverable. Where disclaimers of warranties are not allowed in full or in part, this disclaimer may not apply to You.}
    \item[\textbf{b.}] \textbf{To the extent possible, in no event will the Licensor be liable to You on any legal theory (including, without limitation, negligence) or otherwise for any direct, special, indirect, incidental, consequential, punitive, exemplary, or other losses, costs, expenses, or damages arising out of this Public License or use of the Licensed Material, even if the Licensor has been advised of the possibility of such losses, costs, expenses, or damages. Where a limitation of liability is not allowed in full or in part, this limitation may not apply to You.}
    \item[c.] The disclaimer of warranties and limitation of liability provided above shall be interpreted in a manner that, to the extent possible, most closely approximates an absolute disclaimer and waiver of all liability.
\end{itemize}

\begin{center}
    \textbf{Section 6 -- Term and Termination.}
\end{center}

\begin{itemize}
    \item[a.] This Public License applies for the term of the Copyright and Similar Rights licensed here. However, if You fail to comply with this Public License, then Your rights under this Public License terminate automatically.
    \item[b.] Where Your right to use the Licensed Material has terminated under Section 6(a), it reinstates:
    \begin{itemize}
        \item[1.] automatically as of the date the violation is cured, provided it is cured within 30 days of Your discovery of the violation; or
        \item[2.] upon express reinstatement by the Licensor.
    \end{itemize}

     For the avoidance of doubt, this Section 6(b) does not affect any right the Licensor may have to seek remedies for Your violations of this Public License.

    \item[c.] For the avoidance of doubt, the Licensor may also offer the Licensed Material under separate terms or conditions or stop distributing the Licensed Material at any time; however, doing so will not terminate this Public License.
    \item[d.] Sections 1, 5, 6, 7, and 8 survive termination of this Public License.
\end{itemize}

\begin{center}
    \textbf{Section 7 -- Other Terms and Conditions.}
\end{center}

\begin{itemize}
    \item[a.] The Licensor shall not be bound by any additional or different terms or conditions communicated by You unless expressly agreed.
    \item[b.] Any arrangements, understandings, or agreements regarding the Licensed Material not stated herein are separate from and independent of the terms and conditions of this Public License.
\end{itemize}

\begin{center}
    \textbf{Section 8 -- Interpretation.}
\end{center}

\begin{itemize}
    \item[a.] For the avoidance of doubt, this Public License does not, and shall not be interpreted to, reduce, limit, restrict, or impose conditions on any use of the Licensed Material that could lawfully be made without permission under this Public License.
    \item[b.] To the extent possible, if any provision of this Public License is deemed unenforceable, it shall be automatically reformed to the minimum extent necessary to make it enforceable. If the provision cannot be reformed, it shall be severed from this Public License without affecting the enforceability of the remaining terms and conditions.
    \item[c.] No term or condition of this Public License will be waived and no failure to comply consented to unless expressly agreed to by the Licensor.
    \item[d.] Nothing in this Public License constitutes or may be interpreted as a limitation upon, or waiver of, any privileges and immunities that apply to the Licensor or You, including from the legal processes of any jurisdiction or authority.
\end{itemize}

\end{document}
